\begin{document}
\section{Motivation}
As high-throughput technologies generate increasingly large and complex biological datasets, specifically in the form of multi-omics networks, there is a growing need for scalable methods to aid in the mechanistic interpretation of these data. Traditionally, human experts have taken a deep-but-narrow approach to interpretation, focusing on a small number of genes or pathways in great detail. Conversely, general purpose language models tend to provide a broad-but-shallow perspective, offering high-level summaries without the depth of insight required for mechanistic understanding. Both approaches have their limitations: human interpretation is time-consuming and constrained by cognitive bandwidth, while LLMs often lack deep multi-step reasoning capabilities and are constrained by cost, data quality, and scope. To address this gap, we propose \textsc{Mentor-RL}, a foundation-style model for genome-scale mechanistic reasoning over multiplex networks, trained to provide a base policy that can be adapted to tissue-, organ-, disease-, or cell-type specific contexts. 

\section{Core Approach}
\textsc{Mentor-RL} is designed as a single, turn-based policy with two modes: actor and verifier. The actor mode generates reasoning paths and tool calls, while the verifier mode synthesizes evidence, updates the current hypothesis, and creates a plan for future exploration. The model maintains a structured memory that captures the state of the interpretation, including evidence, provenance, reasoning traces, and tool utilization. This structured memory promotes human interpretability and allows for auditable outputs, enabling users to inspect the evidence and reasoning process behind the model's conclusions. To further ground training in biological reality, \textsc{Mentor-RL} is trained for tool use on multiplex graph objects. This allows experimental and database-derived evidence to shape reasoning. The core unit of reasoning is the biological module, a functionally coherent group of entities within the multiplex. The model learns to interpret these modules through various reasoning tasks, including explanation, recovery, refinement, separation, and recognition of no shared mechanism, while exploring both top-down and bottom-up modules constructed from the multiplex.

\section{Training Data and Environment}
In order to train \textsc{Mentor-RL}, we have developed a training environment and corpora that support multi-step reasoning over biological networks. We use genome-wide \textsc{Mentor} results to construct top-down modules, and we leverage random-walk-with-restart lines-of-evidence (RWR-LOE) to provide complementary, bottom-up modules. With these corpora as our reasoning substrate, we train \textsc{Mentor-RL} to interpret the modules through a variety of reasoning tasks, including explanation, recovery, refinement, and separation. During training, the model learns to recover true modules from corrupted versions without seeing the original. Trajectories are then scored based on the quality of the recovered module and the associated evidence, with rewards provided for correct interpretations and penalties for incorrect ones. The training environment supports the use of external tools to provide reasoning evidence, allowing humans to inspect trajectories for interpretability and auditability. With this training setup, we cover fundamental reasoning tasks for mechanistic interpretation, enabling the model to generalize to a wide range of biological interpretation contexts. 

\section{Optimization and Evaluation}
We utilize a three-stage optimization process to train \textsc{Mentor-RL} for broad coverage of mechanistic reasoning tasks. First, we employ supervised fine-tuning (SFT) to train the model on open- and closed-book biological knowledge, promoting biological world knowledge, schema following, and tool use behavior. Next, we apply direct preference optimization (DPO) to pairs of candidate reasoning extensions, rewarding correct continuations of an interpretation and penalizing incorrect ones. Finally, we apply group-relative policy optimization (GRPO) to optimize multi-step reasoning paths over entire trajectories, promoting consistent long-horizon interpretations of modules over longer, more complex reasoning paths. This optimization process allows us to control different aspects of the model's behavior, such as tool use and general knowledge, single-step decision-making, and long-horizon reasoning. We will evaluate model performance on held-out modules from the training corpora, as well as on previously published expert-verified mechanistic interpretations.

\section{Expected Contributions}
After training \textsc{Mentor-RL}, we expect to make three major contributions:
\begin{enumerate}
    \item an open-weight, foundation-style reasoning model for mechanistic interpretation of multiplex networks, with adaptability to specific contexts through fine-tuning and prompting,
    \item a training environment, harness, and corpora for multi-step mechanistic reasoning over biological networks, and
    \item a deterministic, module-grounded reward framework for training and evaluating long-horizon scientific reasoning without proprietary judges. 
\end{enumerate}

\end{document}
